\scanpage{6}
\mspage{\handwritten{\deleted{\unsure{II|I}} 1}}

\deleted{\handwritten{1}. \uline{Tores maximaux\handwritten{.} \deleted{et groupe de Weyl}}\handwritten{.}}

\deleted{A partir de maintenant, nous ferons usage des résultats du Séminaire Chevalley 56/57 sur les Groupes Algébriques. Suivant l'usage en théorie des groupes algébriques, nous référons à ce Séminaire \handwritten{p}ar le sigle BIBLE .}

\deleted{Soit d'abord $G$ un groupe algébrique sur un corps \uline{algébriquement clos}\deleted{\unsure{,}} $k$. On appelle \uline{tore maximal} de $G$ un \deleted{tore} sous-groupe algébrique $T$ de $G$ qui est un \uline{tore} (ce qui signifie ici, $k$ étant algébriquement clos, qu'il est isomorphe à un groupe de la forme $\underline{G}_m^{\,r}$), et qui est maximal pour cette propriété. Noter que, $k$ étant parfait, $G_{\text{réd}}$ est un sous-groupe de $G$, lisse sur $k$, c'est donc essentiellement un groupe algébrique au sens de BIBLE. D'ailleurs, tout sous-groupe réduit de $G$ (tel un tore) est automatiquement un sous-groupe de $\deleted{\unsure{G|R}}G_{\text{réd}}$, par suite les tores maximaux de $G$ coïncident avec les tores maximaux de $G_{\text{réd}}$. Lorsque $G$ est affine, donc $G_{\text{réd}}$ affine, un théorème fondamental de Borel nous dit que deux tores maximaux de $G$ sont conjugués \deleted{\unsure{sous G(|sous $G$}} par un élément de $G(k)$ \inserted{\handwritten{$=G_{\text{réd}}(k)$}} (BIBLE 6, th.4 c ), en particulier ils ont même dimension. \inserted{Nous appellerons} \deleted{L}\handwritten{l}a dimension commune des tores maximaux de $G$ \deleted{est alors appelée} le \uline{rang réductif} de $G$.}

\margin{\handwritten{$G_{\text{réd}}(k)$}}

\deleted{\deleted{\unsure{(}}Noter d'ailleurs que la restriction $G$ affine est inutile, comme il résulte d'un théorème \deleted{\unsure{de|bien}} connu de Chevalley, affirmant que tout groupe algébrique connexe lisse sur $k$ est une extension d'une variété abélienne par un groupe affine\deleted{\unsure{,}}; \deleted{nous n'aurons pas besoin de cette \unsure{généralisation}),} \inserted{\handwritten{cf N\textsuperscript{o}6 \unsure{3|1}.}}}

% slanted
\margin{\handwritten{pas à la ligne}}

\deleted{Dans \deleted{\unsure{la suite}} \inserted{\handwritten{les N\textsuperscript{os} 1 à 4}}, nous nous bornerons \deleted{aux \unsure{groupes groupes|groupes-groupes|groupes}} \inserted{\handwritten{le plus souvent aux}} préschémas en groupes affines sur la base.}

% slanted
\margin{\handwritten{à la ligne}}

\deleted{\deleted{$G$ sur $k$ étant comme ci-dessus,} \inserted{\handwritten{Soit $G$ un groupe algébrique \deleted{\illegible{1}} lisse sur $k$,}} et $T$ un tore maximal de $G$, le centralisateur \inserted{$C(T)$} de $T$ dans $G$ sera appelé le}

% archive sheet number: 5
